\documentclass[10pt,a4paper]{article}
\usepackage[utf8]{inputenc}
\usepackage[a4paper,left=3cm,right=2cm]{geometry}
\usepackage{amsmath}
\usepackage[thmmarks, amsmath, thref]{ntheorem}
\usepackage{amsfonts}
\usepackage{amssymb}
\usepackage{fancyhdr}
\pagestyle{fancy}
\usepackage{anysize}
\usepackage{graphicx}
\usepackage{float}
\usepackage{hyperref}
\usepackage{multirow}
\usepackage{enumitem}
\newtheorem{theorem}{Theorem}
\theoremsymbol{\ensuremath{\blacksquare}}
\newtheorem*{solution}{Solución:}
\usepackage{listings}
\usepackage{xcolor}

% Margenes y formalidades

\textwidth 6.25in % Margen Derecho y ancho del texto
\textheight 8.5in
\oddsidemargin 0in % Margen Izquierdo
\headheight 0in % Largo del texto
\leftmargin 1in
\parindent 0em % Salto entre la indentación (Sangría)https://www.overleaf.com/4519127669sdpptkhvsbfn
\topmargin -0.5in % Margen Superior
\parskip 0.5ex % Salto entre parrafos
\baselineskip 1pt

%\lhead{\includegraphics[scale=0.2]{logo.png}} % texto izquierda de la cabecera
%\chead{}                                      % texto centro de la cabecera
%\rhead{\includegraphics[scale=0.2]{FIG/logo-mba-2016.png}} % número de página a la derecha

\renewcommand{\headrulewidth}{0.4pt} % grosor de la línea de la cabecera
\renewcommand{\footrulewidth}{0.4pt} % grosor de la línea del pie

\begin{document}



\pagebreak
\vspace*{4mm} 

\begin{center}
\begin{huge}
\textbf{\\Ejercicios Econometría 2}
\end{Large}\\
%{\large 2023-2}\\
\end{center}

\begin{enumerate}

  \item  Se cuenta con la siguiente información para 5 individuos, donde $Y$ es la variable dependiente y $X_1$, $X_2$ son variables explicativas:

\[
\begin{array}{c|ccc}
\hline
\text{Individuo} & X_1 & X_2 & Y \\
\hline
1 & 4 & 1 & 5 \\
2 & 3 & 2 & 7 \\
3 & 5 & 3 & 10 \\
4 & 7 & 4 & 13 \\
5 & 6 & 4 & 12 \\
\hline
\end{array}
\]

\begin{enumerate}
  \item Estime el modelo $Y = \beta_0 + \beta_1\,X_1 + \varepsilon$, ignorando la variable $X_2$. Calcule los coeficientes, el $R^2$, el $R^2$ ajustado y las sumas de cuadrados (SST, SSR y SSE).
  \item Discuta la posibilidad de sesgo al omitir $X_2$.
  \item Estime el modelo $Y = \beta_0 + \beta_1 X_1 + \beta_2 X_2 + \varepsilon$, calcule el nuevo $R^2$ (y las sumas de cuadrados) y compare con el modelo simple. Hint: 
  \[
(X^TX)^{-1}
= \frac{1}{95}
\begin{bmatrix}
281 & -72 & 35\\
-72 & 34 & -35\\
35 & -35 & 50
\end{bmatrix}.
\]
\end{enumerate}






    \item Utilizando una base de datos de 4.137 alumnos universitarios, se estimó la ecuación siguiente de MCO:

    
$$\hat{colgpa} =1,392 -0,0135 hsperc + 0,0014 sat$$

Con $R^2=0,273$ y donde $colgpa$ es el promedio de las calificaciones de los estudiantes universitarios,
$hsperc$ es el percentil en la clase de educación media en la que se graduó (definida de manera que, por
ejemplo, $hsperc= 5$ significa el 5\% superior de su clase), y $sat$ son los puntajes combinados en
matemáticas y habilidades verbales en la prueba de admisión universitaria de los alumnos.
\begin{enumerate} 
\item (8 puntos) ¿Por qué es lógico que el coeficiente en $hsperc$ sea negativo?


\item  (8 puntos) ¿Cual es el promedio de calificaciones universitarias predicho cuando $hsperc= 20$ y $sat=900$?




\item (10 puntos) Suponga que dos estudiantes, A y B, se graduaron en el mismo percentil de la enseñanza media, pero el puntaje $sat$ del estudiante A es 140 puntos más alto. ¿Cuál es la diferencia predicha en el promedio
de calificaciones universitarias para estos dos alumnos? 



\item (10 puntos) Determine el $R^2$ ajustado ($\Bar{R^2}$) de la regresión. Compárelo con el $R^2$ y explique.



\item (20 puntos) Considerando que:

\[(X^\top X)^{-1} =
\begin{bmatrix}
8,5 & 0,19 & -0,11 \\
0,19 & 0,5 & -0,02 \\
-0,11 & -0,02 & 0,04 \\
\end{bmatrix}
\]

$$ \sum_{i=1}^{n} (y_i - \bar{y})^2=6.800$$
Determine la varianza de los estimadores de la regresión MCO.


\end{enumerate}


\item (Continuación ejercicio 1) Se tienen los siguientes datos de salario, años de educación y experiencia laboral de 5 individuos:
\begin{center}
\begin{tabular}{|c|c|c|c|}
\hline
\textbf{Individuo} & \textbf{Educ} ($X_1$) & \textbf{Exp} ($X_2$)  & \textbf{Salario ($Y$)}\\
\hline
1 & 12 & 14 & 1000 \\
2 & 14 & 15 & 1200 \\
3 & 15 & 17 & 1300 \\
4 & 16 & 15 & 1200 \\
5 & 18 & 16 & 1500 \\
\hline
\end{tabular}
\end{center}

Se hizo el modelo en torno solo a eduación de la forma
$$\hat{Salario}= 115 + 75\cdot educ$$
con $R^2 = 0.85$ y $R^2_{adj} = 0.8$
\begin{enumerate}

    \item (5 puntos) ¿Qué sesgos puede tener el modelo presentado?


    \item (10 puntos) Luego se realiza un modelo incluyendo la variable Experiencia, por lo que el salario estaría explicado por la educación y la experiencia de la siguiente manera:
    $$\widehat{Salario} = -430 + 60 \cdot Educ + 50 \cdot Exp$$
    Calcule SST, SSE, SSR, $R^2$ y $R^2$ ajustado.
    
    \item (5 puntos) Compare ambos modelos. 
    
\end{enumerate}

\end{enumerate}



\end{enumerate}

\section{Formulario}


\begin{itemize}

\item  Estimadores MCO para RLS.
{
$$\hat{\beta}_0 = \bar{y} - \hat{\beta}_1 \bar{x} $$
$$\hat{\beta}_1 = \frac{\sum_{i=1}^{n}(x_i - \bar{x})(y_i-\bar{y})}{\sum_{i=1}^{n}(x_i - \bar{x})^2}=\frac{Cov(x,y)}{Var(x)}$$}


\item  Estimadores MCO para RLM.
{
$${\hat{\beta}} = (X^TX)^{-1} X^TY $$}


    \item \[ R^2  = 1 - \frac{SSR}{SST} \] 
\item  $$SST = \sum_{i=1}^{n} (y_i - \bar{y})^2 \quad \quad SSR = \sum_{i=1}^{n}  \hat{u}_i^2 $$
    \item \[ \bar{R}^2 =  1 - \frac{(n-1)SSR}{(n-k-1)SST} \]

    \item  $$\hat{\sigma}^2 = \frac{\sum_{i=1}^{n} \hat{u}_i^2 }{n-k-1}= \frac{SSR}{n - k - 1}$$
    \item \[
\text{Var}(\hat{\beta}|X) = \hat{\sigma}^2(X^{\top}X)^{-1}
\]
    \end{itemize}



\end{document}